% ============================================================================
% Резюме
% ============================================================================
\section*{Резюме}

Настоящата дипломна работа представя проектирането, реализацията и експерименталната оценка на Coroute --- високопроизводителна C++ библиотека за изграждане на уеб приложения, базирана на съвременните възможности на стандарта C++20. Основната цел на изследването е да се демонстрира, че комбинацията от stackless корутини, операционно-специфични асинхронни I/O механизми и алгоритмично оптимизирано маршрутизиране може да постигне производителност, съизмерима с водещите индустриални решения, при значително опростен програмен модел.

Архитектурата на Coroute се основава на няколко ключови иновации. Първо, изпълнителният модел използва типа \texttt{Task<T>}, който реализира lazy корутини с поддръжка на detached изпълнение и кооперативно прекратяване. Този модел елиминира традиционните недостатъци на callback-базираното асинхронно програмиране, като запазва ефективността на неблокиращия I/O. Второ, мрежовият слой абстрахира платформено-специфичните механизми --- IOCP за Windows и io\_uring за Linux --- зад унифициран интерфейс, позволявайки zero-copy операции и batching на системни извиквания. Трето, маршрутизирането използва DFA-базиран алгоритъм, публикуван в „Matching Text from Start to Finish Against Multiple Regular Expressions" \cite{stankov2024regex}, който постига O(N) времева сложност спрямо дължината на URL-а, независимо от броя на регистрираните маршрути. Четвърто, архитектурата включва шаблонна система за изгледи (Template Views) и безпроблемна интеграция с Flutter чрез Dart FFI, което позволява преизползването на библиотеката като headless web engine за изграждане на нейтив мобилни и десктоп приложения.

Библиотеката поддържа множество протоколи върху един порт чрез ALPN негоциация: HTTP/1.1 с keep-alive, HTTP/2 с HPACK компресия и stream мултиплексиране, и WebSocket за двупосочна комуникация в реално време. Типово-безопасното извличане на параметри използва C++20 concepts за compile-time валидация, елиминирайки цял клас runtime грешки. Шаблонната система за изгледи разширява възможностите за сървърно рендиране чрез интеграция с шаблонизатора Inja.

Експерименталната оценка, проведена на система с AMD Ryzen 5 3600 (6 ядра, 12 нишки) и до 1024 конкурентни клиента, демонстрира throughput от 1,378,578 заявки в секунда за прости HTTP отговори, с медианна латентност от 212 микросекунди и 0\% error rate. Сравнителният анализ с Crow, Boost.Beast и Drogon показва, че Coroute постига сравнима или по-добра производителност при значително по-прост API.

Научният принос на работата включва: интеграция на DFA-базирано маршрутизиране в пълноценен уеб библиотека; корутинен изпълнителен модел, специално проектиран за мрежови приложения; и демонстрация на практическата приложимост на подхода чрез production-scale примерно приложение.

\vspace{1cm}

\textbf{Ключови думи:} C++20, корутини, уеб сървър, HTTP/2, WebSocket, IOCP, io\_uring, DFA маршрутизиране, типова безопасност, асинхронен I/O, Flutter, Dart FFI, Template Views, Inja
